\section{Introduction to Quantum Computing - 量子计算初探}
In classical computation, \impt{bits} (0 or 1) are used to represent information through whether or not there is current flowing through a transistor. Transistors are classical objects. \\
In quantum computation, \impt{qubits} are used to represent information. This may take many forms: atomic qubits, like electrons (ground state to be 0 and excited state to be 1); photonic qubits, where the position of a photon on the \textit{dual rail} decides it representing 0 or 1. Other forms include super-conducting circuits or artificial atoms where the quantum property of currents is leveraged. A good qubit is expected to interact properly with the experimenter (e.g. super-conducting circuits), and not interact with the environment (e.g. atoms and photons). \\
Classical computations have different approaches. The most common one is \impt{deterministic classical computation}, where the output is deterministic for a given input. Consider a piece of $2$-bit information, $01$, we can encode this using a $2^2$-dimension vector:
$$\begin{pmatrix}
    0 & 1 & 0 & 0
\end{pmatrix}^\top$$
where the each element represents a $2$-bit state -- $00, 01, 10, 11$ -- respectively. The $1$ for $01$ shows that this is a \impt{deterministic state}. \\
Another approach is \impt{randomized computation}, which can be na\"ively interpreted as the combination of deterministic computation and the ability to flip coins. Consider we decide a $3$-bit state by flipping $2$ fair coins for the first two bits, and set the last bit to be $0$, then the $3$-bit state can be represented by a $2^3$-dimension vector:
$$\begin{pmatrix}
    \frac{1}{4} & 0 & \frac{1}{4} & 0 & \frac{1}{4} & 0 & \frac{1}{4} & 0
\end{pmatrix}^\top$$
where each element represents a $3$-bit state -- $000, 001, \dots, 111$ -- respectively, and the number for each element represents the \impt{probability} of the state being in that state. \\
Notice that an operation like XOR physically operates on $n$-bits, but abstractly and analytically, it acts on a $2^n \times 1$-dimension vector. We use this mathematical (linear algebraic) idea extensively. \\
Contrary to classical computations, quantum computations relies on quantum states, where a state of $n$ qubits is described as a vector of length $2^n$ consisting \impt{complex numbers}. The origin of complex number and negative numbers in quantum states is the single-particle double-slit experiment that shows interference of a single particle with itself. \\
\bigskip \\
We have discussed the general contrast between classical computation and quantum computation, and the different approaches in classical computations. Now we wish to introduce our motivation to develop quantum computation methods. \\
There are \impt{certain} problems, which quantum computers, using \impt{some} quantum algorithm, can solve faster than classical computers, using \impt{any} classical algorithms (not mathematically proven). \\
Some common quantum algorithms include:
\begin{enumerate}
    \item Satisfiability (SAT): best classical algorithm -- $\order{2^n}$, best quantum algorithm (Grover's) -- $\order{\sqrt{2^n}}$. There is a \impt{polynomial speedup}.
    \item Factoring problem: best classical algorithm -- $2^{\order{n^\frac{1}{3}}}$, best quantum algorithm (Shor's) -- $\order{n^3}$, where $n$ is the number of input digits in binary. There is a \impt{super-polynomial speedup}.
    \item Simulating quantum systems: best classical algorithm -- $t \cdot 2^n$, best quantum algorithm -- $t \cdot \text{poly}(n)$, where $t$ is the time it takes to evolve from the initial state to the final state. One such system of interest is the covalent bond in a hydrogen molecule.
\end{enumerate}
As of right now, we have scaled up to about $100$ \impt{useful} qubits, and about $1000$ $2$-qubit gates. IBM projects in 2029, we can have $100$ million gates on $200$ \impt{logical} qubits (each logical qubit consists of about $100$ to $1000$ normal qubits). In very recent developments, Google announced its $105$-qubit Willow quantum processor with a $2$-qubit gate error rate of $0.33\%$ (meaning they have about $\frac{1}{0.33\%}$ gates). For comparison purposes, the benchmark for factoring RSA-2048 requires about $20$ million qubits and $3$ billion gates; the benchmark for simulating a compound like FeMoco requires about $4$ million qubits and $5$ billion operations.

\newpage